\documentclass[10pt,a4paper]{article}
\usepackage[a4paper,margin=20mm,top=19mm,bottom=20mm]{geometry}
\usepackage[T1]{fontenc}
\usepackage[utf8]{inputenc}
\usepackage{lmodern}
\usepackage{microtype}
\usepackage{array,longtable,booktabs,tabularx}
\usepackage{xcolor}
\usepackage{enumitem}
\usepackage{ragged2e}
\usepackage{needspace}
\usepackage{fancyhdr}
\usepackage{hyperref}
\definecolor{navy}{HTML}{24456C}
\definecolor{primary}{HTML}{24456C}
\definecolor{lightgray}{HTML}{F2F4F7}
\definecolor{muted}{HTML}{66788B}
\definecolor{rulegray}{HTML}{DAE1E8}
\hypersetup{colorlinks=true,linkcolor=navy,urlcolor=navy}
\setlength{\parindent}{0pt}
\setlength{\parskip}{3pt}\linespread{0.96}
\setlength{\emergencystretch}{3em}
\renewcommand{\arraystretch}{1.14}
\setlength{\tabcolsep}{2.5pt}
\setlist[itemize]{leftmargin=1.4em,itemsep=2pt,topsep=3pt}
\setlist[enumerate]{leftmargin=1.55em,itemsep=4pt,topsep=3pt}
\newcolumntype{L}[1]{>{\RaggedRight\arraybackslash}p{#1}}
\newcolumntype{R}[1]{>{\RaggedLeft\arraybackslash}p{#1}}
\newcommand{\bigsec}[1]{\par\needspace{4\baselineskip}\vspace{9pt}\noindent{\Large\sffamily\bfseries\color{navy}#1}\par\vspace{2pt}\hrule height 0.4pt\color{black}\vspace{5pt}}
\newcommand{\medsec}[1]{\par\needspace{3\baselineskip}\vspace{6pt}\noindent{\large\sffamily\bfseries\color{navy}#1}\par\vspace{3pt}}
\newcommand{\smallsec}[1]{\par\needspace{2\baselineskip}\vspace{4pt}\noindent{\normalsize\sffamily\bfseries#1}\par\vspace{2pt}}
\pagestyle{fancy}
\fancyhf{}
\renewcommand{\headrulewidth}{0pt}
\fancyfoot[L]{\scriptsize\color{muted}V2X-PQC}
\fancyfoot[R]{\scriptsize\color{muted}\thepage}
\begin{document}

\begin{center}

{\Large\bfseries Slovak University of Technology in Bratislava}\\[2pt]

{\large Faculty of Informatics and Information Technologies}\\[15pt]

{\LARGE\bfseries\color{primary} Methods of Engineering Work}\\[4pt]

{\large Academic Year 2026/2027}

\end{center}

\vspace{10pt}

\begin{center}

\colorbox{lightgray}{\parbox{0.92\linewidth}{\vspace{7pt}\textbf{Team:} Mat\'u\v{s} \v{C}apo, Oleksii Demydchenko, Jakub Dufek, Danylo Dusheiko\vspace{7pt}}}

\end{center}

\vspace{17pt}

\begin{center}

{\Large\bfseries Application for Applied Research (Call VV-2024)}\\[7pt]

{\large\color{primary}\bfseries V2X-PQC: Post-Quantum Cryptography for Secure V2X Communication}

\end{center}

\vspace{12pt}
\bigsec{1. Excellence}
\medsec{1.1 Project Intent, Proposed Approach and Application Level of Expected Results}
The project's intent is to determine which PQC signature scheme, if any, meets the real constraints of V2X and, where none does, to design one that does. We will:\par

\begin{enumerate}
\item extend the open-source AWSIM simulator with a V2X communication layer and a message-verification module using the OMNeT++ framework;
\item use it to run a reproducible benchmark of candidate PQC signature schemes under realistic traffic;
\item use the identified bottlenecks to propose and evaluate an original PQC algorithm or an optimisation of an existing one.
\end{enumerate}
The expected results will reach TRL 3 following testing and experimental validation in a controlled simulation environment. Vendors and standardisation bodies can then use them as evidence for follow-up development.\par

\medsec{1.2 Relevance of the Goals and Originality}
\textbf{Relevance.} The standardised PQC signatures fit V2X poorly for different reasons:\par

\begin{itemize}
\item ML-DSA (Dilithium): ML-DSA-44 signatures of 2,420 bytes are many times larger than ECDSA signatures of about 64 bytes. In dense traffic, this can saturate a channel with limited capacity (e.g. a 10 MHz ITS-G5 channel).
\item Falcon (FN-DSA): signatures are much smaller, at about 666 bytes (Falcon-512). However, its practical suitability for resource-limited verifying devices at busy intersections, under the 100 ms latency budget required by SAE J 2945-1-2026, remains unresolved.
\end{itemize}
\textbf{Originality.} Existing PQC evaluations mostly use isolated micro-benchmarks and rarely reflect realistic traffic. Here, cryptographic cost is evaluated jointly with traffic density, message rates (up to 1,000 messages per second at the receiver) and the limited processing capacity of verifying devices, all within one reproducible simulation environment.\par

\medsec{1.3 State of Knowledge, Scientific Foundations and Prior Results}
\begin{itemize}
\item \textbf{Standards.} NIST has standardised ML-DSA (FIPS 204) and SLH-DSA (FIPS 205); Falcon is being standardised as FN-DSA. V2X security today follows IEEE 1609.2 and ETSI ITS security, with certificate-based ECDSA signing.
\item \textbf{Literature.} Existing studies mainly report single-operation timings of PQC schemes on generic or embedded CPUs, and some evaluate PQC in V2X in analytical or small-scale settings. The gap is a systematic comparison under realistic traffic loads and receiver-side verification limits.
\item \textbf{Tools.} We build on the open-source ecosystem: AWSIM and OMNeT++.
\end{itemize}
\needspace{12\baselineskip}
\medsec{1.4 Objectives and Feasibility}
\par\medskip\noindent\small\begin{tabular}{@{}L{11.53mm}L{70.80mm}L{82.33mm}@{}}
\toprule
\textbf{\#} & \textbf{Objective} & \textbf{Measurable target} \\
\midrule
O1 & Simulation framework: AWSIM augmented with a V2X communication layer using OMNeT++ tools & Reproducible, documented, open-source release; standardised BSMs, adjustable density, message rates of up to 1,000 msg/s \\
O2 & Benchmark of candidate PQC signatures (including ML-DSA and Falcon, with ECDSA as the baseline) & Signature size, channel load, verification latency and throughput reported for each scheme and traffic scenario \\
O3 & Original PQC algorithm, or optimisation of an existing one & Verification latency below 100 ms at 1,000 msg/s on the reference device; substantially less transmitted cryptographic data than ML-DSA \\
O4 & Dissemination and guidance & At least one research paper; evidence-based recommendations for vendors, operators and standardisation bodies \\
\bottomrule\end{tabular}\normalsize\par\medskip
\textbf{Feasibility.} O1 and O2 rely on established tools and well-defined methods, so we assess their risks as low. O3 is the high-risk, high-gain part, because a new signature scheme needs rigorous security analysis and cannot be guaranteed. The project therefore has a staged, risk-managed design: O3 builds on the O2 findings and includes a fallback, namely an optimised implementation or protocol-level improvement of an existing scheme with a security argument. This fallback still produces publishable and usable results.\par

\medsec{1.5 Methodology and Its Justification}
\begin{enumerate}
\item \textbf{Requirements and scenarios.} Define V2X constraints (channel capacity, message size limits, latency budget, message rates) and traffic scenarios (urban intersection, highway, dense congestion) from standards and literature.
\item \textbf{Framework development.} Add a V2X layer to AWSIM, covering message generation (BSM), channel and loss models, receiver queueing, and a verification module with plug-in cryptographic back-ends. Validate the layer against analytical estimates and published channel behaviour.
\item \textbf{Benchmark.} Run each candidate scheme across the scenario matrix (density $\times$ message rate $\times$ device capacity), with repetitions and fixed seeds for reproducibility. Measure cryptographic micro-benchmarks on representative embedded or automotive-class hardware and feed the measured costs into the simulation, so the results rest on measurements and not only on modelled timings.
\item \textbf{Bottleneck analysis.} Identify whether bandwidth, verification time, certificate and public-key size, or queueing dominates for each scheme.
\item \textbf{Algorithm design and evaluation.} Candidate directions include optimised arithmetic (e.g. NTT and vectorised or hardware-assisted implementations), batch or pipelined verification, certificate and public-key handling strategies, parameter tailoring, or a new scheme variant. We will pursue the most promising directions according to the bottleneck analysis and evaluate them with the same framework. Any new construction will undergo a security analysis (reduction arguments and cryptanalysis review).
\item \textbf{Dissemination.} Publish the framework, data and results to the extent permitted.
\end{enumerate}
\textbf{Justification.} Simulation allows controlled, repeatable experiments across traffic densities that are impractical or unsafe to reproduce on real roads. Combining it with hardware-measured cryptographic costs limits a main weakness of simulation, namely model inaccuracy. Using open-source simulators helps others replicate and extend the work.\par

\clearpage
\bigsec{2. Impact}
\medsec{Benefits for the Development of Knowledge, Applied Research and Innovation}
The project addresses an unresolved gap in post-quantum cryptography (PQC) for cooperative intelligent transport systems. Neither of the standardised candidates fits V2X conditions: ML-DSA (Dilithium) produces 2,420-byte signatures that overload communication channels at high traffic density, while Falcon, with 666-byte signatures, makes verification too demanding for roadside units (RSUs) at busy intersections and risks exceeding the 100 ms latency limit required by the SAE J 2945-1-2026 standard. The project will deliver:\par

\begin{itemize}
\item \textbf{New knowledge:} A systematic, reproducible evaluation of PQC signature schemes under realistic V2X load (message rates, hardware limits, latency budgets), which is currently lacking.
\item \textbf{Simulation environment:} By combining AWSIM, a Unity-based simulator for Autoware that generates realistic vehicle dynamics, sensor data (LiDAR, cameras and GNSS), and traffic scenarios, with OMNeT++, a discrete-event network simulator that models wireless channels, packet loss, latency and congestion, we will build an environment for testing and comparing cryptographic algorithms in a realistic V2X setting.
\item \textbf{New PQC algorithm:} A PQC algorithm for RSUs and on-board units (OBUs), compatible with IEEE 1609.2, aiming to reduce transmitted cryptographic data by at least 40\% compared with ML-DSA-44 while ensuring that signature verification latency remains below 100 ms at 1,000 messages per second. The designed PQC algorithm would secure V2X communication against future quantum attacks while respecting the real-time limits of hardware.
\end{itemize}
\medsec{Usability of the Results for the Applicant and Users}
The primary users are manufacturers and integrators of RSUs and OBUs, road infrastructure operators, and automotive suppliers. These organizations will integrate the research outputs into their roadside and vehicle security products and verify them on target hardware, which provides a direct route from research to practice.\par

Because the solution targets IEEE 1609.2, it can be reused wherever V2X systems must migrate to PQC. Quantum-safe security matters especially for vehicles and roadside equipment because their service lives are long, so systems deployed now must remain secure for many years.\par

\medsec{Degree of Economic Benefit}
\begin{itemize}
\item \textbf{Added value and competitiveness:} The automotive supply chain gains knowledge, a testing environment and an optimised PQC algorithm, shortening the time to market for quantum-safe products.
\item \textbf{Quality of life and health:} Secure V2X communication is a precondition for safety applications such as collision warnings, so protecting message authenticity supports road safety. Fewer hazardous situations mean fewer injuries and fatalities, along with reduced long-term disability and strain on healthcare systems. Smoother, more predictable traffic flow enabled by trustworthy cooperative driving also cuts congestion, emissions, and driver stress, improving everyday wellbeing in urban areas.
\end{itemize}
\clearpage
\bigsec{3. Implementation}
\medsec{3.1 Project Plan and Expected Outputs}
The V2X-PQC project will run for 24 months, from January 2027 to December 2028. It is divided into five work packages (WPs), following the schedule in VV-D. Each stage builds on the results of the previous one.\par

\smallsec{WP1 --- Analysis of V2X Security Requirements and PQC Algorithms}
January 2027 -- April 2027\par

The team will review V2X security requirements, relevant standards and existing research on post-quantum cryptography. ECDSA, ML-DSA-44 and Falcon-512 will be selected for comparison. The team will also define the experimental scenarios and performance metrics.\par

\noindent\textbf{Milestone M1:} Research review and testing methodology completed.\par
\smallsec{WP2 --- Development of the V2X Simulation Environment}
May 2027 -- October 2027\par

The team will integrate AWSIM and OMNeT++ to simulate data generated by vehicles and their on board sensors and network communication. A separate component will be added to generate and verify digital signatures. Development will begin with simple scenarios before moving to more complex traffic conditions.\par

\noindent\textbf{Milestone M2:} Functional simulation environment with message verification.\par
\smallsec{WP3 --- Comparative Testing of Signature Algorithms}
November 2027 -- April 2028\par

The selected algorithms will be tested under different simulated traffic conditions, including higher message rates approaching 1,000 messages per second at a receiving unit. The team will compare signature sizes, signing and verification times, processing throughput, network delays and message delivery.\par

Experiments will be repeated under consistent conditions to obtain reliable results.\par

\noindent\textbf{Milestone M3:} Comparative testing completed and main performance limitations identified.\par
\smallsec{WP4 --- Proposal of a New or Optimised PQC Approach}
May 2028 -- September 2028\par

Based on the test results, the team will identify the main limitations of the evaluated algorithms and propose a possible improvement for V2X communication.\par

The proposal may focus on an existing algorithm or a different way of handling cryptographic operations and information. The team will explain its expected benefits, security considerations and limitations. A complete implementation of a new cryptographic algorithm is an expected output of the research effort.\par

\noindent\textbf{Milestone M4:} Proposed PQC approach and its technical justification completed.\par
\smallsec{WP5 --- Final Evaluation and Publication}
October 2028 -- December 2028\par

The final stage will focus on reviewing the results, preparing documentation and writing a research paper. The team will summarize the main findings and identify possible directions for future research.\par

\noindent\textbf{Milestone M5:} Final report, documentation and research manuscript completed.\par
\medsec{3.2 Project Management}
The project will be carried out by four team members, each responsible for a specific area.\par

\par\medskip\noindent\small\begin{tabular}{@{}L{37.87mm}L{42.81mm}L{83.98mm}@{}}
\toprule
\textbf{Team member} & \textbf{Position} & \textbf{Main responsibility} \\
\midrule
Matúš Čapo & Project Coordinator & Planning, coordination and progress monitoring \\
Jakub Dufek & Cryptography Researcher & Analysis and testing of digital signature algorithms \\
Oleksii Demydchenko & V2X Developer & Development and integration of simulation software \\
Danylo Dusheiko & Testing and Evaluation & Test scenarios, measurements and analysis of results \\
\bottomrule\end{tabular}\normalsize\par\medskip
The team will hold regular meetings to discuss progress, solve problems and plan upcoming tasks. GitHub will be used to manage source code and Jira to track development tasks.\par

At the end of each work package, the team will review the completed work and check whether the planned milestone has been achieved.\par

The total planned workload is 2,800 hours over two years, with 1,400 hours allocated to each year.\par

\needspace{13\baselineskip}
\medsec{3.3 Implementation Risks and Risk Mitigation}
The main risks are related to software integration, experimental measurements and project scheduling.\par

\par\medskip\noindent\small\begin{tabular}{@{}L{49.40mm}L{41.16mm}L{74.10mm}@{}}
\toprule
\textbf{Risk} & \textbf{Affected output} & \textbf{Mitigation} \\
\midrule
Difficult integration of AWSIM and OMNeT++ & Simulation environment (WP2) & Start with simple scenarios. If direct synchronisation is too difficult, use recorded vehicle movement data. \\
Problems with cryptographic libraries & Benchmark results (WP3) & Use existing implementations and test them separately before integration. \\
High computing requirements or inconsistent measurements & Comparative study (WP3) & Begin with smaller experiments, repeat tests and document the configurations. Dedicate a larger portion of the budget to computation hardware. \\
No clear improvement identified & PQC proposal (WP4) & Focus on the observed limitations and propose realistic improvements supported by the results. \\
Delays in development & Final outputs (WP5) & Monitor progress regularly and adjust individual tasks when needed. \\
\bottomrule\end{tabular}\normalsize\par\medskip
The team will test components step by step to identify problems early. If a planned approach proves too difficult, an alternative method will be considered and its limitations documented.\par

\medsec{3.4 Budget Justification}
The total project budget is 68,638.40 EUR for 2027--2028, as detailed in VV-C.\par

Most of the budget is allocated to personnel costs, including 43,200 EUR for wages and 15,638.40 EUR for estimated social and health insurance contributions. These costs reflect the work needed for research, simulation development, cryptographic testing and evaluation. The remaining funds cover travel, materials, computing services, energy, communications and indirect costs. These expenses will support software development, experiments and the presentation of research results.\par

The project will mainly use open-source software and existing university computing resources where available, helping to reduce equipment costs. The budget is therefore focused on the main research activities and the resources needed to complete the planned work packages.\par

\medsec{3.5 Technical and Personnel Infrastructure}
The project is planned to use university workstations for software development, simulations and performance testing. Access to these resources will be confirmed before the relevant project stages begin. Additional computing capacity or storage may be obtained through the planned services budget if needed.\par

The main software tools will include AWSIM, OMNeT++ and existing cryptographic libraries. These tools will provide the basis for developing the simulation environment and comparing digital signature algorithms.\par

The four-member team will cover project coordination, cryptographic research, V2X development and testing. The members will cooperate during software integration and experimental evaluation.\par

The project will focus on software-based experiments without requiring physical V2X communication hardware.\par

\end{document}